\section{Two below the degree: a new all-padding late-C extension}
The written assembly, thirteen scoped unbounded Lean kernels, one
complete finite certificate and independent reconstruction have passed
the checks recorded in the accompanying verification manifest. This is
a new refinement; none of the older frozen arguments is edited.
General C/A2/original UI remain OPEN.
Write $q=\phi-1$, $\psi=-q$, and $[t]=1+\cdots+z^{t-1}$.
Let $L=\prod_{i=1}^r[t_i]$, $t_i\ge3$, $r\ge2$,
$h=\deg L=u+2$, $W=h+r$, and $E=(1+z)^uL$.
The target is every late C comparison for $1^\alpha2^uV$, every
$\alpha\ge1$, at every current rank $N\ge W+1$.

\paragraph{An eightfold robust head, with no numerator degree bound.}
For $u\ge16$, let $Q_u$ denote the degree-$u$ truncation of
$(1+z)^u/(1+qz)$. For ANY nonnegative LC $N$ with $N_0=1$,
$N_1\le u/2+1$, put $B=NQ_u$. The proposed strengthening is
\[
 M_B(s)\ge8q^s\{B_{s+1}+2qB_s+q^2B_{s-1}\},
                         \qquad 2\le s\le u-1.       \tag{1}
\]
Use exactly the registered first-intermediate-row Cauchy--Binet
expansion. For $1\le k\le s$, its kernel $K(s,k)$ is at least
$Q_sQ_k/(s+1)$, and $T_k/Q_k\le u+q$, including $k=u-1$.
At $s=2$ the two cases $k=1,2$ have positive polynomial
expansions in $u-16$ recorded by the new audit. At $s=3$ use
$Q_s\ge\binom{u-1}{3}\ge8(u+q)$ and $32q^3\le8$.
For $4\le s\le u-3$, use $8(s+1)q^s\le40q^4\le6$ and
$Q_s\ge\binom{u-1}{2}\ge6(u+q)$.
At $s=u-2\ge14$, $8(s+1)q^s\le120q^{14}\le3/16$ and
$Q_s\ge u-q\ge(3/16)(u+q)$.
For the LAST row, the exact base and decreasing ratio give
\[
       u(u+1)q^{u-1}\le1/5\quad(u\ge16).
\]
The base is $16\cdot17q^{15}<1/5$, and its step ratio is
$q(u+2)/u\le3/4$. Also $Q_{u-1}=\phi(1-q^{2u})\ge8/5$,
using $\phi\ge161/100$ and $q^{32}\le1/1000$.
Consequently $8u q^{u-1}(u+q)\le8/5\le Q_{u-1}$.

Row zero requires $Q_s\ge q^s(12u+8+8q)$. For $s\le u-3$
use $q^s\le2/5$ and the quadratic binomial bound; at $u-2$
use $q^{u-2}\le1/16$ and $u-q$; at $u-1$ use
$12u+8+8q\le u(u+1)$. This pays the omitted $N_{s+1}$
because $N_{s+1}\le N_1N_s$, proving (1).

\paragraph{The quotient can have a negative constant coefficient.}
Put
\[
 c=\psi^uL(-\phi)>0,\qquad J=(E-c)/(1+qz).
\]
The norm bound gives $0<c\le4q^2<8/5$, so $J_0=1-c$ need
not be positive. The proposed auxiliary theorem, for $u\ge16$, is:
$J_n>0$ for every $n\ge1$; the rank-one truncation of $J$ is
globally LC; $J$ is increasing initially and then unimodal;
and the ENTIRE positive primitive $g=(\Delta J)^+$ is globally LC.
If $J_0<0$, $g$ has positive support beginning at rank one.

Delete a LARGEST interval $[t]$ and write $L=L_0[t]$,
$r_0=r-1\ge1$. Its old degree is $h_0=u+3-t\le u$.
The registered degree-boundary/surplus theorems make the OLD quotient
positive global LC with positive primitive LC after its first centre;
the single old exception has $u=3$ and cannot occur here.
The quotient correction has degree $t-2$, so ordinary centres
$n\ge t$ and primitive centres $n\ge t+1$ are inherited from
the corresponding old uniform window.

For the ordinary head, (1) applies with $N=L$ and
$L_1=r\le u/2+1$. The perturbation $c\psi^n$ has
$c/q\le4q<8$. All head entries at ranks $1..u$ are positive
since $c q^n\le4q^3<1$ and $(LQ_u)_n\ge1$ there.
This proves every ordinary centre $2..u-1$.
The first ordinary centre is positive if $J_0<0$; if $J_0\ge0$,
its exact expression $w^2-(1-c)(E_2+qw)>0$, $w=u+r$,
uses $E_2\le w^2/2$. Rank zero is omitted in the claimed
nonnegative LC sequence. Only $t=u+1$ leaves an ordinary gap at
centre $u$; the explicit endpoint computation below pays it.

For the primitive, set $B=L_0Q_u$, $a=\phi^2c$, $b=B_1=u+r_0-q$,
and $X_n=B_n-a\psi^n$, using the formal quotient coefficients
through the displayed reads. Then
\[
 d_n=X_n\ (1\le n<t),\qquad d_t=X_t-1,\qquad
 d_{t+1}=X_{t+1}-b.                                \tag{2}
\]
Here $0<a\le4$ and $a/q\le4\phi<8$, so (1) pays
every $X$ minor at centres $2..u-1$.
All its required head reads are positive; use $u-q$ through
rank $u-2$, $Q_{u-1}>3/2$ at rank $u-1$ and the explicit
long endpoint if rank $u$ is read. The omitted atom at these
last ranks has magnitude less than $1/100$.
Lowering the next value at centre $t-1$ improves its minor.

For $t\le u-1$, ordinary LC of $B$ and
$b^2-2B_2=u-r_0-q^2$ show
\[
 bB_{t-1}-2B_t\ge3u/2.
\]
Indeed $B_{t-1}\ge(r_0+1)(u-q)$ and the concave function
\[
 (u-r-1)(r+1)(u-1)-(3/2)u(u+r)
\]
is nonnegative on $1\le r\le u/2$ at $u\ge16$.
Its endpoints are $(u^2-15u+8)/2$ and
$(u^3-10u^2-4u+4)/4$.
The perturbation coefficient $a q^{t-1}$ is at most $4/5$:
if $t\ge6$, use $a\le4$ and $4q^5<2/5$; if $t\le5$,
$u+2\le4r$ forces $r\ge5$, hence
$a q^{t-1}\le\phi^2(2q^2)^5=32q^8\le4/5$.
Its error is at most $(4/5)(b+2q)\le6u/5+8/15$.
Thus $bX_{t-1}-2X_t\ge3u/10-8/15>0$.
The exact junction minor is $M_X(t)+1+bX_{t-1}-2X_t$,
so it is positive before clipping.

If $t=u+1$, necessarily $L_0=[3]$. With $v=\psi^u$,
\[
 d_{u-1}=\phi\binom u2-u+2\phi+2\phi^2v,\quad
 d_u=\phi u-1-2\phi v,\quad d_{u+1}=q+2v,\quad
 d_{u+2}=-u-1+q-2qv<0.
\]
The estimates $|v|\le1/100$, $d_u\ge3u/2$,
$d_{u-1}\le u^2$, $d_{u+1}\le1$ pay the last primitive
centre $u$. They also give $d_u>d_{u+1}$, hence ordinary
concavity, and LC, at the sole missing ordinary centre $u$.
If $t=u$, necessarily $L_0=[4]$, and
\[
 d_u=\phi\binom u2-u+2\phi-1-(2\phi-1)v,\quad
 d_{u+1}=qu-2+q+(3-\phi)v,\quad
 d_{u+2}=\phi-(3\phi-4)v-\binom u2-q^2(u+2)<0.
\]
Here $d_u\ge2u^2/3$, $d_{u-1}\le u^3/3$, and
$d_{u+1}\le3u/4$ pay the final junction. Every later primitive
centre is inherited or is an endpoint with next positive part zero.
At the first primitive centre, if $J_0<0$ its previous positive
part is zero. Otherwise its exact minor is
\[
 w^2-E_2-\phi(w-1)+c\{E_2+\phi(w-2)\}+\phi c^2>0,
\]
since $w\ge18$. This proves the auxiliary assertions.

\paragraph{The one negative initial coefficient does not invalidate the late cut.}
Let $x_s=\sum_{j\ge1}q^jJ_{s+j}$ for $s\ge0$.
Convolution with the positive rank-one truncation of $J$ proves LC
of $x$ at every centre $s\ge1$. Moreover
\[
 x_0=C_u+1,\quad x_s=C_{u+s}\ (s\ge1),\quad
 C=L\sum_{k=0}^{u-1}(z+z^2)^k.
\]
Thus $x$ is EXACTLY the infinite-initial-ones ACTUAL signed
derivative tail from $W$, including its extra unit at $W$.
For $a=\alpha-1$, the finite actual tail is
$A_s=(x_s-E_{s-a})^+$, $s\ge0$.

The old signed-primitive kernel identity is algebraic and still holds
for signed $J$. If $J_0<0$, the additional negative primitive
at rank zero contributes only through output rank $a$, because
$K_a$ vanishes above $a$. Therefore at every centre $s\ge a+2$
all three reads ignore it. The positive primitive $g$ is globally LC;
the later negative primitive contributes one geometric term whose
sign follows from $J_k\le J_M$ after its first mode $M$:
\[
 \sum_{k>M}q^k\Delta J_k
       =q^2\sum_{k>M}q^kJ_k-q^{M+1}J_M\le0.
\]
The registered geometric-tilt/constant-cut proof consequently applies
unchanged to these triples. The true $c$ atom affects only rank $a$.
The first three startup centres $a-1,a,a+1$ have the same exact
payments as the registered all-padding proof: $x_s/x_{s-1}<\phi$,
$w=E_1=u+r\ge18$, and
$w(w-1)/2\le E_2\le w^2$. The first finite tail mode is no later
than the infinite one by the registered actual source-density order,
and the latter is no later than $M$. Hence every required late
pre-mode minor is within the paid range.

For $u\le13$, the registered all-small-binary crossing theorem already
gives the same actual derivative-tail LC. For $u=14,15$ there are
exactly $54+65=119$ multifactor degree-$u+2$ profiles. The new
complete padding-clamped certificate is to pay precisely these
bounded exceptions. Single-factor tails use the old one-jump theorem.

\paragraph{A new actual entrance for C, retaining BOTH count losses.}
Let $R$ be the infinite-binary reference with the same CI suffix,
$x=\Delta R_W$, $y=\Delta R_{W+1}$, $w=u+r$.
The FIRST loss from the finite binary count is at $W$ and has
signed-derivative series
\[
 z^W E/(1-z-z^2).
\]
The separate finite-initial-ones loss is $z^{W+a}E$.
Consequently the ACTUAL losses at $W,W+1$ are
\[
 (\delta_0,\delta_1)=
 (2,2w+1)\ (\alpha=1),\quad
 (1,w+2)\ (\alpha=2),\quad
 (1,w+1)\ (\alpha\ge3).
\]
Every earlier ordinary/cumulative coefficient is unchanged.
The inverse identity gives $D_{W+1}/D_W<\phi$ whenever both
are positive: at infinite padding its strict gap is $J_1>0$;
$\alpha=1$ changes that gap by $w-\phi>0$ and
$\alpha=2$ changes it by one.

The old reference identities give $R_{W-1}=y$, $R_W=x+y$,
$y\ge x>0$, and $A=x^2+xy-y^2>0$. The exact ordinary seed is
\[
 F_WD_W-F_{W-1}D_{W+1}
   =A-2\delta_0x+(\delta_1-\delta_0)y+\delta_0^2>0.
\]
The displayed losses and $w\ge4$ prove the last inequality.
All cumulative ratios through $W-1$ are reference ratios, hence
strictly exceed $\phi$. The derivative-tail LC makes its positive
ratios nonincreasing. The ordinary seed and TWO successive cumulative
integrations now give every C read, at every shift at least three,
starting at $N=W+1$. Decreasing derivative tails are automatic.
For the few smaller $u$ the same entrance algebra applies with
$w\ge4$; the already registered $u\le8$ late-C theorem can also
be used directly. General earlier-rank C, A2 and original UI are
not inferred from this one-step extension.

\paragraph{Hereditary consequences, without a new C inheritance claim.}
The registered arbitrary-insertion signed-tail identity now propagates
late pre-mode derivative LC from ANY submultiset of degree $u+2$,
as well as the older degree-$u+1$ and long-factor seeds.
In particular, for EVERY $u\ge0$ and EVERY number $r$ of threes
with $2r>u$, the actual positive derivative tail is pre-mode LC at
every centre $N\ge3r+1$, at every padding. Choose the least even
degree above $u$: it is $u+1$ when $u$ is odd and $u+2$ when
$u$ is even. The one-factor or zero-binary cases use old analytic
bases; all further threes are arbitrary insertions.

More generally an unresolved deletion-minimal crossing seed can now
be assumed to have
\[
 3\le d_i\le u,\quad u+3\le h\le2u,\quad
 h-\min d_i\le u,\qquad d_i=t_i-1.
\]
Indeed, if its least degree were two, deletion-minimality would force
$h\le u+2$, which is already closed. Thus the entire minimal crossing
entrance problem contains NO factor $[3]$. One may also exclude seeds
having a subset of degree $u+1$ or $u+2$. This is a derivative-LC
reduction; early F/G and C do NOT become permutation invariant.
